% Purpose: Supplementary material for the High-Performance Computing
% Implementation chapter presenting
% architectural hardware specifications, NUMA topology models, Slurm job
% orchestration schemas,
% theoretical parallel scaling models (Amdahl, Gustafson, Isoefficiency),
% Lustre parallel I/O contention proofs, and Roofline arithmetic intensity
% derivations for deep learning phase picking.
% Status: Complete; publication-grade reference text.
% Source-of-truth: doc/UNITS/main/Chapter5_HPC_Implementation.tex,
% OGS/utils/Leonardo/Makefile, OGS/utils/Leonardo/LAUNCHME.sh,
% OGS/utils/Leonardo/ACTIVATEME.sh, and OGS/utils/Leonardo/ktanakah.sh.

\section*{HPC Implementation: Hardware Architecture, Slurm
Orchestration, and Parallel Scaling Models}
\addcontentsline{toc}{section}{HPC Implementation: Hardware
Architecture, Slurm Orchestration, and Parallel Scaling Models}
\label{sec:supp_hpc_implementation}

This supplementary document provides formal architectural models,
Slurm workload
configuration schemas, and theoretical parallel performance
derivations supporting
\cref{chap:hpc}.

% =============================================================================
% 1. ARCHITECTURAL HARDWARE DECOMPOSITION AND NUMA TOPOLOGY
% =============================================================================

\subsection*{Leonardo Booster Node Architecture and NUMA Memory Topology}
\label{subsec:supp_numa_topology}

The Leonardo Booster module, engineered by Atos on the BullSequana
X2135 Da Vinci platform,
delivers dense hybrid acceleration tailored for tensor-parallel deep
neural network
inference and high-throughput data conditioning.

\subsubsection*{Processor and Memory Subsystem Hierarchy}
Each Booster blade integrates a single socket host processor and
four discrete accelerators:
\begin{itemize}
  \item \textbf{Host CPU}: Intel Xeon Platinum 8358 (Ice Lake
    microarchitecture) featuring
    \num{32} physical cores (\num{64} execution threads via
    Hyper-Threading), a base clock
    of $\qty{2.60}{\giga\hertz}$ (all-core Turbo
    $\qty{3.40}{\giga\hertz}$), and $\qty{54}{\mega\byte}$
    of shared L3 cache ($\qty{1.6875}{\mega\byte}$ non-inclusive per core).
  \item \textbf{Host RAM}: $\qty{512}{\giga\byte}$ registered
    DDR4-3200 ECC memory across
    \num{8} independent channels ($8 \times \qty{64}{\giga\byte}$
    dual-rank DIMMs),
    providing a peak theoretical host memory bandwidth of:
    \begin{equation}
      B_{\mathrm{host\_peak}} \coloneqq \num{8} \times
      \qty{3200e6}{\transfer\per\second} \times \qty{8}{\byte}
      = \qty{204.8}{\giga\byte\per\second}.
      \label{eq:supp_host_bandwidth}
    \end{equation}
  \item \textbf{Accelerators}: Four NVIDIA Ampere A100-SXM4-64GB
    \acp{GPU}, each equipped with
    \num{108} Streaming Multiprocessors (SMs), \num{6912} FP32 CUDA
    cores, and \num{432}
    third-generation Tensor Cores.
  \item \textbf{High-Bandwidth Memory (HBM2e)}:
    $\qty{64}{\giga\byte}$ of stacked HBM2e
    memory per \ac{GPU} accessed across a \num{5120}-bit wide memory
    bus operating at
    $\qty{1.51}{\giga\hertz}$, yielding an individual device bandwidth of:
    \begin{equation}
      B_{\mathrm{GPU\_peak}} \coloneqq \qty{1.555}{\tera\byte\per\second},
      \label{eq:supp_gpu_bandwidth}
    \end{equation}
    and an aggregate node accelerator bandwidth of $4 \times
    \qty{1.555}{\tera\byte\per\second}
    \equiv \qty{6.22}{\tera\byte\per\second}$.
  \item \textbf{NVLink 3.0 Interconnect}: Each A100 \ac{GPU} exposes
    \num{12} NVLink 3.0
    links, each providing $\qty{50}{\giga\byte\per\second}$
    bidirectional bandwidth,
    establishing an all-to-all crossbar fabric providing
    $\qty{600}{\giga\byte\per\second}$
    bidirectional bandwidth between any accelerator pair.
\end{itemize}

\subsubsection*{NUMA Topology and UPI Interconnect Penalty}
The Intel Xeon 8358 processor partitions its \num{32} physical cores
and \num{8} memory
controllers into two internal Non-Uniform Memory Access (\ac{NUMA})
domains, commonly referred
to as Sub-NUMA Clustering (SNC-2):
\begin{align}
  \mathcal{N}_0 &\coloneqq \{ \text{Cores } 0 \dots 15, \quad
  \text{Memory Channels } 0 \dots 3, \quad \text{PCIe Root Complexes } 0, 1 \},
  \label{eq:supp_numa0} \\
  \mathcal{N}_1 &\coloneqq \{ \text{Cores } 16 \dots 31, \quad
  \text{Memory Channels } 4 \dots 7, \quad \text{PCIe Root Complexes } 2, 3 \}.
  \label{eq:supp_numa1}
\end{align}

Accelerators $\text{GPU}_0$ and $\text{GPU}_1$ are attached to the
PCIe switches routed to Root Complex
$\mathcal{N}_0$, whereas $\text{GPU}_2$ and $\text{GPU}_3$ are
attached to Root Complex $\mathcal{N}_1$.
When a worker thread running on core $c \in \mathcal{N}_0$ stages
waveform data into $\text{GPU}_2 \in \mathcal{N}_1$,
the Direct Memory Access (DMA) transaction must cross the internal
Ultra Path Interconnect (UPI) bus.

\begin{theorem}[NUMA UPI Bus Penalty on Host-to-Device Transfers]
  \label{thm:supp_numa_penalty}
  Let $T_{\mathrm{local}}$ and $T_{\mathrm{remote}}$ denote the DMA
  transfer latency for
  a contiguous host buffer of size $S_{\mathrm{buf}}$ to device memory.
  Cross-socket transfers incur an additive latency penalty
  $\tau_{\mathrm{UPI}}$ and
  a multiplicative bandwidth reduction factor $\eta_{\mathrm{UPI}} < 1$:
  \begin{align}
    T_{\mathrm{local}}(S_{\mathrm{buf}}) &\coloneqq
    \tau_{\mathrm{PCIe}} + \frac{S_{\mathrm{buf}}}{B_{\mathrm{PCIe\_peak}}},
    \label{eq:supp_t_local} \\
    T_{\mathrm{remote}}(S_{\mathrm{buf}}) &\coloneqq
    (\tau_{\mathrm{PCIe}} + \tau_{\mathrm{UPI}}) +
    \frac{S_{\mathrm{buf}}}{\eta_{\mathrm{UPI}} B_{\mathrm{PCIe\_peak}}},
    \label{eq:supp_t_remote}
  \end{align}
  where $\tau_{\mathrm{PCIe}} \approx \qty{1.2}{\micro\second}$,
  $\tau_{\mathrm{UPI}} \approx \qty{0.8}{\micro\second}$,
  and measured empirical efficiency $\eta_{\mathrm{UPI}} \approx \num{0.60}$.
\end{theorem}

\begin{proof}
  On PCIe Gen4 $\times 16$, the simplex bandwidth is
  $B_{\mathrm{PCIe\_peak}} \coloneqq \qty{31.5}{\giga\byte\per\second}$.
  When data originates in $\mathcal{N}_0$ and targets
  $\mathcal{N}_1$, data packets must be serialized
  across the internal UPI mesh, which arbitrates between coherent
  CPU cache invalidations and
  asynchronous I/O streams. The UPI link limits sustained
  host-to-device throughput to
  $\approx \qty{18.9}{\giga\byte\per\second}$, corresponding to
  $\eta_{\mathrm{UPI}} \approx \num{0.60}$,
  representing a $\qty{40}{\percent}$ bandwidth degradation.
\end{proof}

To strictly enforce NUMA locality and eliminate the UPI penalty, our
submission launcher
pins each worker task to its local socket and adjacent GPU:
\begin{equation}
  \operatorname{Task}_k \implies \left\{ \text{Cores: } [8k, 8k+7],
  \quad \text{Device: } \text{cuda:}k \right\}, \quad \forall k \in
  \{0, 1, 2, 3\}.
  \label{eq:supp_pinning_rule}
\end{equation}

% =============================================================================
% 2. SLURM BATCH SCHEDULING SPECIFICATIONS AND WRAPPER MECHANICS
% =============================================================================

\subsection*{Slurm Batch Scheduling Specifications and Dynamic
Substitution Engine}
\label{subsec:supp_slurm_specifications}

\Cref{tab:supp_slurm_params} lists the authoritative operational
Slurm parameters configured
for ISCRA-C production workflows across CINECA Leonardo.

\begin{table}[htbp]
  \centering
  \small
  \caption{Authoritative Slurm workload allocation parameters for
    Leonardo Booster and DCGP modules.
    Derived from production batch scripts
    \texttt{OGS/utils/Leonardo/LAUNCHME.sh} and
  \texttt{OGS/utils/Leonardo/ktanakah.sh}.}
  \label{tab:supp_slurm_params}
  \begin{threeparttable}
    \begin{tabularx}{\linewidth}{l l l >{\raggedright\arraybackslash}X}
      \toprule
      \textbf{Directive} & \textbf{Booster ML Production} &
      \textbf{DCGP CPU Inversion} & \textbf{Operational Function} \\
      \midrule
      \texttt{--account} & \texttt{IscrC\_AISeism} &
      \texttt{IscrC\_AISeism\_0} & Resource allocation accounting tag \\
      \texttt{--partition} & \texttt{boost\_usr\_prod} &
      \texttt{dcgp\_usr\_prod} & Target hardware module partition \\
      \texttt{--qos} & \texttt{boost\_qos\_bprod} &
      \texttt{dcgp\_qos\_bprod} & Quality of Service tier and priority \\
      \texttt{--nodes} & \num{1} to \num{8} & \num{1} to \num{4} &
      Allocated physical blade count \\
      \texttt{--tasks-per-node} & \num{4} & \num{1} to \num{16} &
      Dispatched task/rank count per host \\
      \texttt{--cpus-per-task} & \num{8} & \num{7} to \num{112} &
      Dedicated CPU threads per worker process \\
      \texttt{--gres} & \texttt{gpu:4} & \texttt{tmpfs:3T} & Generic
      resource requests (GPUs or fast scratch) \\
      \texttt{--mem} & \qty{490}{\giga\byte} & \qty{490}{\giga\byte}
      & Enforced host memory allocation ceiling \\
      \texttt{--time} & \texttt{24:00:00} & \texttt{24:00:00} &
      Maximum permissible job wall-clock time \\
      \bottomrule
    \end{tabularx}
    \begin{tablenotes}[flushleft]
      \footnotesize
    \item Target \& Scope: Concrete Slurm execution directives used
      in the production environment.
    \item What the reader should notice: Booster enforces a strict
      8-to-1 core-to-GPU mapping, whereas DCGP maximizes thread
      parallelism for octree ray tracing.
    \end{tablenotes}
  \end{threeparttable}
\end{table}

\subsubsection*{Dynamic Substitution and Trap-Safe Workflow Lifecycle}
The shell orchestrator \texttt{LAUNCHME.sh} utilizes an immutable
template architecture
to generate executable batch jobs dynamically. The workflow
execution follows four strictly
ordered lifecycle phases:

\begin{enumerate}
  \item \textbf{Parameter Ingestion and Validation}: The wrapper
    validates environment variables,
    verifies project directory structure on Lustre, checks the
    existence of daily waveform
    SQLite indexes, and parses command-line arguments (year, days,
    picking model, associator).
  \item \textbf{Atomic In-Place Template Substitution}: Using a
    stream editor (\texttt{sed}),
    placeholders in \texttt{ktanakah.sh} are substituted in place:
    \begin{equation}
      \text{Placeholder } \texttt{\#SBATCH --nodes=\#}
      \xrightarrow{\quad\text{\texttt{sed}}\quad} \texttt{\#SBATCH
      --nodes=} N_{\mathrm{nodes}}.
    \end{equation}
  \item \textbf{Signal Trapping for Template Restoration}: To
    prevent corruption of the base
    template if the submission is interrupted, a POSIX exit handler
    is established:
    \begin{lstlisting}[language=bash,basicstyle=\scriptsize\ttfamily]
trap 'sed -i "s/--nodes=${NODES}/--nodes=#/" ktanakah.sh' EXIT INT TERM
    \end{lstlisting}
    This guarantees that \texttt{ktanakah.sh} returns to its clean
    immutable state regardless of
    exit status.
  \item \textbf{Environmental Toolchain Initialization}: The
    dispatched job executes
    \texttt{ACTIVATEME.sh}, purging default user environment
    modules, loading vendor-optimized
    compilers and libraries (\texttt{nvhpc}, \texttt{cuda/12.2}),
    and activating the isolated
    Conda environment:
    \begin{lstlisting}[language=bash,basicstyle=\scriptsize\ttfamily]
eval "$(/leonardo_work/IscrC_AISeism/.miniconda3/bin/conda shell.bash hook)"
conda activate SBC_3.12
    \end{lstlisting}
\end{enumerate}

% =============================================================================
% 3. PARALLEL SCALING MODELS: AMDAHL, GUSTAFSON, AND ISOEFFICIENCY
% =============================================================================

\subsection*{Parallel Scaling Models and Scalability Formulations}
\label{subsec:supp_scaling_models}

High-performance seismic processing exhibits two distinct
computational regimes:
(i) embarrassingly parallel temporal data partitioning, and (ii) intra-cluster
multi-core phase association and location.

\subsubsection*{Strong Scaling: Amdahl's Law and Serial Bottlenecks}
For a fixed dataset size $W$ (e.g., a \num{24}-hour swarm across
\num{100} stations),
strong scaling evaluates execution speedup as a function of
allocated processing cores $N$.

Let $T_1$ denote execution time on a single core, decomposed into a
perfectly parallelizable
component $T_{\mathrm{p}}$ and an inherently serial component $T_{\mathrm{s}}$:
\begin{equation}
  T_1 \coloneqq T_{\mathrm{s}} + T_{\mathrm{p}}, \quad f \coloneqq
  \frac{T_{\mathrm{p}}}{T_1} \in [0, 1],
  \label{eq:supp_amdahl_decomposition}
\end{equation}
where $f$ is the parallelizable fraction. Under ideal
parallelization, execution time on $N$ cores is:
\begin{equation}
  T_N \coloneqq T_{\mathrm{s}} + \frac{T_{\mathrm{p}}}{N} = T_1
  \left( (1 - f) + \frac{f}{N} \right).
  \label{eq:supp_t_n}
\end{equation}
The theoretical speedup $S_{\mathrm{strong}}(N)$ is given by Amdahl's Law:
\begin{equation}
  S_{\mathrm{strong}}(N) \coloneqq \frac{T_1}{T_N} = \frac{1}{(1 -
  f) + \frac{f}{N}}.
  \label{eq:supp_amdahl_law}
\end{equation}
Taking the asymptotic limit as processor allocation grows without bound:
\begin{equation}
  \lim_{N \to \infty} S_{\mathrm{strong}}(N) \equiv \frac{1}{1 - f}.
  \label{eq:supp_amdahl_limit}
\end{equation}
If serial disk I/O, SQLite database locks, or metadata parsing
constitute just $\qty{5}{\percent}$
of runtime ($1 - f = \num{0.05}$), the maximum achievable speedup is
strictly capped at
$S_{\max} = \num{20}$, regardless of whether \num{100} or
\num{10000} cores are deployed.

\subsubsection*{Weak Scaling: Gustafson's Scaled Speedup Law}
In continuous seismic observation, researchers do not keep the
workload fixed; instead,
as computational resources expand, they scale the problem size
proportionally by ingesting
more seismic stations, expanding spatial monitoring envelopes, or
evaluating multi-year spans.

Under Gustafson's Law, the total workload executed on $P$ processors
is scaled such that the
parallel execution time remains constant: $T_P = s + p$, where $s
\coloneqq 1 - p$ is the
normalized serial fraction. On a serial processor, this scaled
workload would require time:
\begin{equation}
  T_1(P) \coloneqq s + p \cdot P = (1 - p) + p \cdot P.
  \label{eq:supp_gustafson_t1}
\end{equation}
The scaled speedup is given by:
\begin{equation}
  S_{\mathrm{weak}}(P) \coloneqq \frac{T_1(P)}{T_P} = \frac{(1 - p)
  + p \cdot P}{1} \equiv P - s(P - 1).
  \label{eq:supp_gustafson_law}
\end{equation}
Unlike Amdahl's Law, Gustafson's speedup scales linearly with
processor count $P$.
For calendar-day temporal decomposition across independent cluster
blades, $s \approx \num{0.005}$,
achieving near-ideal weak scaling efficiency:
\begin{equation}
  E_{\mathrm{weak}}(P) \coloneqq \frac{S_{\mathrm{weak}}(P)}{P} = 1
  - \frac{s(P - 1)}{P} \approx \qty{99.5}{\percent}.
  \label{eq:supp_weak_efficiency}
\end{equation}

\subsubsection*{Isoefficiency Metric}
Parallel efficiency is defined as $E(W, P) \coloneqq \frac{S(W,
P)}{P} = \frac{T_1(W)}{P \cdot T_P(W, P)}$.
Parallel overhead $T_o(W, P)$ is the total cumulative idle or
communication time across all processors:
\begin{equation}
  T_o(W, P) \coloneqq P \cdot T_P(W, P) - T_1(W).
  \label{eq:supp_overhead_def}
\end{equation}
The efficiency can be rewritten as:
\begin{equation}
  E(W, P) = \frac{1}{1 + \frac{T_o(W, P)}{T_1(W)}}.
  \label{eq:supp_efficiency_overhead}
\end{equation}
To maintain a constant target efficiency $E_0 \in (0, 1)$, the
problem size $W \propto T_1(W)$ must satisfy
the isoefficiency relation:
\begin{equation}
  W \equiv \left(\frac{E_0}{1 - E_0}\right) T_o(W, P).
  \label{eq:supp_isoefficiency}
\end{equation}
For distributed phase association (where inter-station pick exchange
  requires $\mathcal{O}(P \log P)$
all-to-all communication), the isoefficiency function scales as $W
\propto P \log P$, indicating
excellent architectural scalability.

% =============================================================================
% 4. LUSTRE PARALLEL I/O STRIPING AND DATABASE LOCK CONTENTION
% =============================================================================

\subsection*{Lustre Parallel I/O Optimization and Day-Sharded SQLite Indexing}
\label{subsec:supp_lustre_io}

High-concurrency processing of multi-year seismic archives is
predominantly limited by file system
metadata latency rather than raw floating-point operations.

\subsubsection*{Lustre File Striping Mechanics}
The CINECA Lustre parallel filesystem distributes storage across
Object Storage Targets (OSTs)
managed by Object Storage Servers (OSSs). A directory configured
with striping parameters
$(c, S)$ splits contiguous waveform files into chunks of size $S$
distributed round-robin
across $c$ distinct OSTs:
\begin{lstlisting}[language=bash,basicstyle=\scriptsize\ttfamily]
lfs setstripe -c 4 -S 4M /leonardo_work/IscrC_AISeism/WORK/waveform
\end{lstlisting}
Under this configuration, a \qty{64}{\mega\byte} continuous miniSEED
day file is divided into
\num{16} striping chunks of \qty{4}{\mega\byte} each, striped across
\num{4} OSTs.
The aggregate theoretical I/O throughput for parallel read streams is:
\begin{equation}
  T_{\mathrm{I/O}}(c) \coloneqq \min\left( c \times
    B_{\mathrm{OST}},\, B_{\mathrm{InfiniBand}},\,
  B_{\mathrm{host\_bus}} \right),
  \label{eq:supp_lustre_throughput}
\end{equation}
where single-OST bandwidth $B_{\mathrm{OST}} \approx
\qty{2.5}{\giga\byte\per\second}$.
For $c = \num{4}$, aggregate streaming read throughput reaches
$\approx \qty{10.0}{\giga\byte\per\second}$,
saturating the host InfiniBand interface and preventing I/O stalls
during neural window batching.

\subsubsection*{POSIX Lock Contention: Monolithic vs. Day-Sharded SQLite}
Maintaining a unified seismic metadata index (mapping station codes,
  channel IDs, start times,
and sampling rates to raw file offsets) is essential for fast
waveform querying.
Standard SQLite utilizes POSIX advisory file locking (`fcntl`) to
ensure ACID compliance.

\begin{theorem}[Elimination of Lock Contention via Day-Sharding]
  \label{thm:supp_sqlite_sharding}
  Let $M$ concurrent worker tasks query and update metadata on a
  shared distributed filesystem.
  Under a monolithic database architecture, aggregate lock wait time
  scales quadratically $\mathcal{O}(M^2)$,
  inducing severe process stalling. Under day-sharded partitioning,
  lock wait time vanishes identically:
  $\mathcal{O}(1)$.
\end{theorem}

\begin{proof}
  Let each worker perform an index transaction requiring hold time $\tau_h$.
  In a monolithic database file $\mathcal{D}_{\mathrm{mono}}$, the
  file lock is a shared resource
  with a single exclusive write lock or shared read lock. On
  distributed parallel filesystems
  (Lustre / GPFS), lock arbitration requires round-trip remote
  procedure calls (RPCs) to the
  Metadata Server (MDS) with network latency $\tau_{\mathrm{RPC}}
  \approx \qty{5}{\milli\second}$.
  The probability of lock contention among $M$ independent Poisson
  processes with transaction rate $\lambda$ is:
  \begin{equation}
    P_{\mathrm{contention}} \coloneqq 1 - \exp(-M \lambda \tau_h).
    \label{eq:supp_lock_prob}
  \end{equation}
  When $M \ge \num{32}$ (e.g., across 8 nodes $\times$ 4 GPUs), $M
  \lambda \tau_h \gg 1$, and total
  queuing delay across all processes scales as:
  \begin{equation}
    T_{\mathrm{wait\_mono}} = \sum_{k=1}^{M-1} k \tau_{\mathrm{RPC}}
    = \frac{M(M - 1)}{2}\tau_{\mathrm{RPC}} \in \mathcal{O}(M^2).
    \label{eq:supp_mono_delay}
  \end{equation}
  In contrast, let the database be sharded by calendar day:
  \begin{equation}
    \mathcal{D} \equiv \bigsqcup_{d=1}^{D_{\mathrm{total}}}
    \mathcal{D}_d, \quad \mathcal{D}_d \coloneqq
    \text{\texttt{ogsDB/YYYY/MM/DD/.squirrel}}.
    \label{eq:supp_sharded_db}
  \end{equation}
  Because tasks are assigned non-overlapping temporal intervals
  ($d_i \ne d_j, \forall i \ne j$),
  worker $i$ accesses database file $\mathcal{D}_{d_i}$ while worker
  $j$ accesses $\mathcal{D}_{d_j}$.
  The intersection of their POSIX inode locks is strictly empty:
  \begin{equation}
    \operatorname{Lock}(\mathcal{D}_{d_i}) \cap
    \operatorname{Lock}(\mathcal{D}_{d_j}) \equiv \emptyset, \quad
    \forall i \ne j.
    \label{eq:supp_disjoint_locks}
  \end{equation}
  Consequently, cross-worker lock contention is identically zero
  ($T_{\mathrm{wait\_sharded}} \equiv 0$),
  enabling linear parallel indexing scaling across hundreds of compute nodes.
\end{proof}

% =============================================================================
% 5. ROOFLINE ARITHMETIC INTENSITY DERIVATION FOR DEEP PHASE PICKERS
% =============================================================================

\subsection*{Roofline Arithmetic Intensity Derivations for Deep Phase Pickers}
\label{subsec:supp_roofline}

To establish whether deep learning phase detection (\ac{PhaseNet},
\ac{EQT}) on NVIDIA A100 GPUs
is bound by memory bandwidth or floating-point compute throughput,
we derive the arithmetic
intensity of 1D convolutional layers operating on seismic waveform tensors.

\subsubsection*{Kinematic Input Waveform Representation}
Input seismic data consists of continuous three-component particle
velocity time series
$\mathbf{v}(t) \equiv \dot{\mathbf{u}}(t) =
[\dot{u}_{\mathrm{E}}(t), \dot{u}_{\mathrm{N}}(t),
\dot{u}_{\mathrm{Z}}(t)]^\top$.
A batch of waveform windows is assembled into a 3D tensor:
\begin{equation}
  \mathbf{X} \in \mathbb{R}^{B \times C_{\mathrm{in}} \times L},
  \label{eq:supp_input_tensor}
\end{equation}
where $B \coloneqq B_{\mathrm{size}}$ is batch size,
$C_{\mathrm{in}} \equiv \num{3}$ channels,
and $L \coloneqq \num{6000}$ samples (corresponding to a
\qty{60}{\second} window sampled at $f_s = \qty{100}{\hertz}$).

\subsubsection*{Arithmetic Intensity Formulation}
Consider a standard 1D convolutional layer with $C_{\mathrm{out}}$
filters of kernel width $K$
operating with stride \num{1} and padding to preserve sequence length $L$.
The total floating-point multiply-accumulate (MAC) operations
performed by the layer is:
\begin{equation}
  N_{\mathrm{FLOPs}} \coloneqq 2 \times B \times C_{\mathrm{out}}
  \times C_{\mathrm{in}} \times K \times L.
  \label{eq:supp_conv_flops}
\end{equation}
The minimum memory traffic (in bytes) transferred across the GPU
memory bus under ideal cache reuse is:
\begin{equation}
  N_{\mathrm{bytes}} \coloneqq \text{sizeof}(\mathrm{float}) \times
  \left[ B \cdot C_{\mathrm{in}} \cdot L + B \cdot C_{\mathrm{out}}
  \cdot L + C_{\mathrm{out}} \cdot C_{\mathrm{in}} \cdot K \right].
  \label{eq:supp_conv_bytes}
\end{equation}
Assuming 32-bit floating point precision
($\text{sizeof}(\mathrm{float}) = \qty{4}{\byte}$),
the theoretical arithmetic intensity $\mathcal{I}_{\mathrm{conv}}$
(in FLOPs per byte) is:
\begin{align}
  \mathcal{I}_{\mathrm{conv}}
  &\coloneqq \frac{N_{\mathrm{FLOPs}}}{N_{\mathrm{bytes}}}
  \label{eq:supp_arith_intensity_def} \\
  &= \frac{2 B C_{\mathrm{in}} C_{\mathrm{out}} K L}{4 \left( B L
      (C_{\mathrm{in}} + C_{\mathrm{out}}) + C_{\mathrm{in}}
  C_{\mathrm{out}} K \right)}.
  \label{eq:supp_arith_intensity_fraction}
\end{align}
For typical convolutional layers where sequence length $B L \gg K$:
\begin{equation}
  \mathcal{I}_{\mathrm{conv}} \approx \frac{2 B C_{\mathrm{in}}
  C_{\mathrm{out}} K L}{4 B L (C_{\mathrm{in}} + C_{\mathrm{out}})}
  = \frac{K}{2} \left( \frac{C_{\mathrm{in}}
  C_{\mathrm{out}}}{C_{\mathrm{in}} + C_{\mathrm{out}}} \right).
  \label{eq:supp_arith_intensity_approx}
\end{equation}

\subsubsection*{A100 Roofline Analysis and Critical Ridge Point}
On the NVIDIA A100-SXM4 GPU operating in TF32 precision:
\begin{itemize}
  \item Peak Tensor Core Throughput: $P_{\mathrm{peak}} \coloneqq
    \qty{156}{\tera\text{FLOPS}} \equiv \qty{156e12}{\text{FLOPs}\per\second}$.
  \item Peak Memory Bandwidth: $B_{\mathrm{peak}} \coloneqq
    \qty{1.555}{\tera\byte\per\second} \equiv
    \qty{1.555e12}{\byte\per\second}$.
\end{itemize}
The critical Roofline ridge point $\mathcal{I}^*$ separating
memory-bound from compute-bound regimes is:
\begin{equation}
  \mathcal{I}^* \coloneqq
  \frac{P_{\mathrm{peak}}}{B_{\mathrm{peak}}} =
  \frac{\qty{156e12}{\text{FLOPs}\per\second}}{\qty{1.555e12}{\byte\per\second}}
  \approx \qty{100.3}{\text{FLOPs}\per\byte}.
  \label{eq:supp_ridge_point}
\end{equation}

\begin{remark}[Operational Regime of Deep Phase Pickers]
  For early convolutional layers ($C_{\mathrm{in}} = 3,
  C_{\mathrm{out}} = 16, K = 7$):
  \begin{equation}
    \mathcal{I}_{\mathrm{early}} \approx \frac{7}{2} \left( \frac{3
    \times 16}{3 + 16} \right) = \frac{7}{2} \left( \frac{48}{19}
    \right) \approx \qty{8.84}{\text{FLOPs}\per\byte} \ll \mathcal{I}^*.
  \end{equation}
  Early layers are strictly memory-bandwidth bound.
  However, for intermediate and deep residual blocks
  ($C_{\mathrm{in}} = 128, C_{\mathrm{out}} = 128, K = 7$):
  \begin{equation}
    \mathcal{I}_{\mathrm{deep}} \approx \frac{7}{2} \left( \frac{128
    \times 128}{128 + 128} \right) = \frac{7}{2} (64) =
    \qty{224.0}{\text{FLOPs}\per\byte} > \mathcal{I}^*.
  \end{equation}
  Deep layers operate strictly in the compute-bound regime of the
  Roofline model.
  By configuring dynamic batch sizes $B_{\mathrm{size}} \ge
  \num{256}$, host memory staging overhead is amortized,
  enabling the A100 Tensor Cores to sustain $> \qty{85}{\percent}$
  of theoretical peak throughput during continuous
  seismic inference.
\end{remark}
